\scanpage{4}
5) Compactification des types modulaires $M_{g,\nu}$ et structure à l'infini. Il semble qu'on ait
\[
\pi_1(\widehat M_{g,\nu})=1\qquad\text{si }g\geq2
\]
(si $g=1$, $\nu=0$, le $\pi_1$ doit être cyclique fini d'ordre pair divisant 12...), et $\pi_1(M_{g,\nu}^{\infty})=1$ itou.

$*$ Déf. de $\widehat M_{g,\nu}$ : schémas \inserted{$X/S$} propres (\unsure{projectifs} ?) \inserted{$X/S$ muni d'\unsure{sections} $(f_1,-f_\nu)$,} et plats sur $S$, de prés. finie, à fibres $X_s$ géom connexes réduites, n'ayant que des pts doubles ordinaires, et les $f_{is}$ étant des pts lisses, enfin le groupe des autom. de $(X_s,f_{1s},\ldots,f_{\nu s})$ étant fini (ce qui signifie que \deleted{\unsure{les}} composante irréductible de genre $0$ de $X_s$ a au moins trois points marqués -- soit par intersection avec d'autres comp. irréd, soit comme pts $f_{is}$ ...)

$*$ Question : si $\nu=0$, $X_s^{\inserted{\text{lisse}}}$ est-il contenu dans sa jacobienne, \unsure{\`a} extension d'une VA par un tore ?

$*$ La vérification du critère valuatif semble étrange.

Les générateurs ``twist'' de $\unsure{\Gamma}_{g\nu}$ sont les lacets autour des comp. irréductibles de $M_{g,\nu}^{\infty}$

\mspage{3}
